\documentclass[12pt]{article}
\usepackage[utf8]{inputenc}
\usepackage{amsmath, amssymb, amsthm}
\usepackage{geometry}

\geometry{a4paper, margin=1in}

\newtheorem{theorem}{Theorem}
\newtheorem{lemma}[theorem]{Lemma}
\newtheorem{proposition}[theorem]{Proposition}
\theoremstyle{definition}
\newtheorem{definition}[theorem]{Definition}

\title{Project Euler Problem 14: Longest Collatz Sequence}
\author{}
\date{}

\begin{document}
\maketitle

\section{Problem Statement}

Define the Collatz map $C : \mathbb{Z}^+ \to \mathbb{Z}^+$ by
\[
C(n) = \begin{cases} n/2 & \text{if } 2 \mid n, \\ 3n + 1 & \text{if } 2 \nmid n. \end{cases}
\]
The chain length $L(n)$ is the number of terms in the sequence $(n, C(n), C^2(n), \ldots)$ up to and including the first~$1$. Find $\arg\max_{1 \le n < 10^6} L(n)$.

\section{Formal Development}

\begin{definition}
The \emph{Collatz orbit} of $n$ is $(C^{(k)}(n))_{k \ge 0}$ where $C^{(0)}(n) = n$ and $C^{(k)}(n) = C(C^{(k-1)}(n))$. The \emph{chain length} is $L(n) = \min\{k \ge 0 : C^{(k)}(n) = 1\} + 1$.
\end{definition}

\begin{theorem}[Recursive Structure]
\label{thm:recur}
$L(1) = 1$, and for $n > 1$: $L(n) = 1 + L(C(n))$.
\end{theorem}

\begin{proof}
Removing the first term $n$ from the orbit $(n, C(n), \ldots, 1)$ yields the orbit from $C(n)$, which has $L(C(n))$ terms.
\end{proof}

\begin{lemma}[Odd-Step Compression]
\label{lem:odd}
For odd $n > 1$: $C(C(n)) = \frac{3n+1}{2}$ and $L(n) = 2 + L\!\left(\frac{3n+1}{2}\right)$.
\end{lemma}

\begin{proof}
If $n$ is odd, $C(n) = 3n+1$ is even (since $3n$ is odd), so $C(3n+1) = \frac{3n+1}{2}$. Applying Theorem~\ref{thm:recur} twice: $L(n) = 1 + L(3n+1) = 2 + L\!\left(\frac{3n+1}{2}\right)$.
\end{proof}

\begin{theorem}[Termination]
The Collatz conjecture ($L(n) < \infty$ for all $n \ge 1$) has been verified computationally for all $n < 2^{68}$, which far exceeds $10^6$.
\end{theorem}

\begin{theorem}[Correctness of Memoized Search]
\label{thm:memo}
Processing starting values $n = 2, 3, \ldots, N-1$ in order, with $\text{cache}[1] = 1$, and computing $L(n) = k + \text{cache}[m]$ where $m$ is the first iterate with a known cache entry and $k$ is the number of steps from $n$ to $m$, correctly fills the cache.
\end{theorem}

\begin{proof}
By Theorem~\ref{thm:recur}, $L(n) = k + L(C^{(k)}(n))$ for any $k$. When $m = C^{(k)}(n)$ satisfies $\text{cache}[m] = L(m)$, we have $L(n) = k + L(m)$. Since we process values in increasing order and every orbit eventually reaches 1 (by verified termination), the base case $\text{cache}[1] = 1$ suffices.
\end{proof}

\begin{proposition}[Complexity]
The algorithm uses $\Theta(N)$ space. The total time is $O(N \cdot \bar{L})$ in the worst case, empirically $O(N \log N)$ with memoization.
\end{proposition}

\begin{proof}
Space: the cache array has $N$ entries. Time: for each starting value, the orbit is followed until a cached value is reached. With memoization, values below $N$ that have been previously computed terminate the chain immediately. The average number of uncached steps per starting value is empirically $O(\log N)$, motivated by the observation that each even step halves the value and each odd-even pair multiplies by $\sim 3/4$.
\end{proof}

\section{Result}

$L(837799) = 525$, which is the maximum over all starting values below $10^6$.

\section{Answer}

\[
\boxed{837799}
\]

\end{document}
